\subsection{Discussion}
\label{Discussion1}
\setlength{\parindent}{20pt}
\justifying

\noindent
This experiment piloted a new statistical learning paradigm to study morpheme learning in an implicit, bilingual setting. A familiarity test demonstrated that participants did indeed pick up some basic knowledge of the morphemes they were exposed to in the initial training phase. An instinctive bias towards familiarity also showed up in the distribution of ``yes'' responses across the testing conditions: as was observed in \textcite{Lelonkiewicz2020}, participants preferred to label as congruent any strings that contained a chunk that they knew already; this effect is amplified when two known chunks are present in the proposed strings. This replication of previous work laid a solid foundation on which to extend that work into the domain of individual variability.

One principal aim of the present research has been a wider inclusion of individual variability than what has been sometimes sought after for the sake of limiting sample noise. The success of our approach was seen in the diversity of backgrounds and languages measured. A PCA done on participants' language scores did not account for as much variance as might have been hoped for, but, given the astonishing number of differences in language learning and habits, it can be said that this uniqueness was partly distilled into a few useful data-driven components. Additionally, the novel cosine similarity metric was shown to be valuable as a ``pure'' measure of language balance that is not affected by the number of languages spoken, unlike the usual entropy metric.

While none of these bottom-up, or indeed none of the top-down measures either, seemed to influence the segmentation abilities of participants, the PCA proved an important tool when examining the pattern of ``yes'' responses to the testing items. The earlier one starts learning a second language, the more prone they are to accept novel stimuli as being linguistic in nature. 

However, some factors also influence a more measured sensitivity to new combinations involving word parts they already know. Age was found to induce a larger morpheme familiarity effect, whereby the mere presence of a known linguistic chunk is enough to suggest to a participant that this novel assembly is congruent with their past knowledge. Participants with more sensitivity to this (namely, older participants) respond ``yes'' at chance level for chunks with no known parts, slightly more to those containing one known part, and even more for those being a combination of two known parts. This age-related enhancement may be due to their longer time manipulating linguistic systems, leading to their having seen a wider variety of combinations of morphemes, and therefore more reliance on known structures. This can also be said about the significant interaction between L2 proficiency and this morpheme familiarity effect. More proficient language users may have a larger vocabulary, and a similar reason for being more sensitive to the presence of familiar word parts.

On a slightly different plane, one further measure pushed the morpheme familiarity effect for certain participants. Those with a smaller difference in scores between their first and second languages also exhibited this pattern of ``yes'' responses. These would be participants who make use of their top two languages more equally, and can be expected to be more flexible in working with both. This flexibility could explain the easier acceptance of new words linked to those previously stored in memory, especially if they're linked to several different part-word items.

However, a solid null result was observed on a group scale of the discrimination performance between within- and between-language morphemes combinations in condition 2M. Crucially, one group show some aptitude for correct sorting of stimuli in the 2M condition. Those having reported using known morphemes in the testing strings seemed better equipped to reject between-language items than those using other strategies or their intuition. These participants were using a morphology-based classification strategy: they may have understood the underlying structure better than participants who focused more on letter-based statistics. A promising avenue for improving the experimental design might be to boost morpheme learning.